\begin{table}[htbp]
\centering
\caption{Top 10 Departamentos por Defección del Cabildo Abierto al Frente Amplio (Modelo con PI)}
\label{tab:dept_top10}
\begin{tabular}{lS[table-format=4]S[table-format=5]S[table-format=2.1]S[table-format=5.0]}
\toprule
{\textbf{Departamento}} & {\textbf{N}} & {\textbf{Votos CA}} & {\textbf{CA$\rightarrow$FA (\%)}} & {\textbf{Votos Est. FA}} \\
\midrule
Río Negro & 121 & 1023 & 95.7 & 979 \\
 & & & {[87.8, 99.5]} & \\
Treinta y Tres & 135 & 794 & 70.2 & 558 \\
 & & & {[39.1, 94.5]} & \\
Colonia & 299 & 1403 & 52.6 & 738 \\
 & & & {[29.4, 77.0]} & \\
Maldonado & 404 & 3059 & 47.9 & 1466 \\
 & & & {[21.0, 75.6]} & \\
Durazno & 152 & 1047 & 40.9 & 429 \\
 & & & {[12.1, 70.4]} & \\
Rocha & 185 & 1237 & 40.0 & 495 \\
 & & & {[9.0, 71.7]} & \\
Florida & 179 & 697 & 38.7 & 270 \\
 & & & {[0.9, 74.9]} & \\
Artigas & 181 & 3163 & 34.4 & 1087 \\
 & & & {[19.5, 48.9]} & \\
San José & 252 & 1405 & 30.8 & 433 \\
 & & & {[6.7, 60.2]} & \\
Paysandú & 281 & 1381 & 27.9 & 385 \\
 & & & {[3.2, 60.8]} & \\
\bottomrule
\end{tabular}
\begin{flushleft}
\small \textit{Nota:} Estimaciones del modelo con Partido Independiente (PI) como categoría separada. Votos estimados calculados como producto de la tasa de defección y votos CA totales. IC, intervalo creíble al 95\%. Las tasas medias posteriores se muestran en la fila superior; los intervalos creíbles en la fila inferior.
\end{flushleft}
\end{table}
